\begin{frame}{第二章}
    \begin{itemize}
        \item 熟悉 $\displaystyle\sum$ 记号（艾弗森括号、条件型下标）
        \item 扰动法
        \item 二重求和：对称式
        \item 熟悉一些化简和式的方法
        \item 下降幂
    \end{itemize}
\end{frame}
\begin{frame}{第二章习题}
    需要写下来的：2 3 4 6 13 14 19 20 21 22 25 32；7 9 10 11 17 23 24 27 29 31；35* 36*

    两周内完成，其中第一周至少要完成第一个分号前的习题（对应教材 2.1-2.5）。另外补充一些 OI 练习题：洛谷 P5686 P4948 P4091 U77198
\end{frame}
\begin{frame}{化简和式的一些指导思想}
    \begin{itemize}
        \item 尽量使用简单的求和条件
        \item 把无关项用分配律提出到求和号外
        \item 交换求和号常常是有用的
    \end{itemize}
\end{frame}
\begin{frame}{含有求和的递归}
    例：$s_n=\sum_{i=1}^n f_i$，$f_n$ 是 Fibonacci 数列

    邻项相减

    本题还有包括但不限于扰动法的各种推导方式
\end{frame}
\begin{frame}{利用对称去掉偏序}
    在三角求和的例子中，我们利用对称性把 $i\le j$ 的求和变成无限制的求和以及 $i=j$ 的求和。这是常用的技巧。例如：

    $$
    \sum_{1\le i\le j\le n}|x_i-x_j|.
    $$

    异或粽子

    另外，这个技巧可以推广为 Burnside 引理。
\end{frame}
\begin{frame}{小练习}
    $$
    \sum_{1\le i_1\le i_2\le\cdots\le i_m\le n} 1
    $$

    是多少？

    我们有很多种方式来推导它：找规律后归纳、组合意义、有限微积分、写递归式，等等
\end{frame}
\begin{frame}{找规律}
    打表发现答案是 $\binom{n+m-1}{m}$。

    证明：按 $n+m$ 归纳：

    $$
    \begin{aligned}
    &\sum_{1\le i_1\le \cdots\le i_m\le n}1\\
    &=\sum_{1\le i_1\le \cdots\le i_m<n}1+\sum_{1\le i_1\le \cdots\le i_{m-1}\le i_m=n}1\\
    &=\binom{(n-1)+m-1}{m}+\binom{n+(m-1)-1}{m-1}\\
    &=\binom{n+m-1}{m}
    \end{aligned}
    $$
\end{frame}
\begin{frame}{组合意义}
    \begin{itemize}
        \item 对每个 $t=1,2,\ldots,n$，考虑有多少个求和指标等于 $t$。这样相当于把 $m$ 个球放到 $n$ 个可空的桶里
        \item 或者，考虑 $[1,i_1],(i_1,i_2],(i_2,i_3],\ldots,(i_{m-1},i_m],(i_m,n]$ 这 $m+1$ 个桶，其中第一个桶非空，其他可空。
        \item 再或者，每个 $1\le i_1\le i_2\le\cdots\le i_m\le n$ 一一对应一条路径 $(1,1)\to (1,i_1)\to (2,i_1)\to (2,i_2)\to (3,i_2)\to (3,i_3)\to\cdots\to (m,i_m)\to (m+1,i_{m})\to (m+1,n)$ 的路径，其中 $(t,i_t)\to (t+1,i_t)$ 是一个向右的步，$(t,x)\to (t,x+1)$ 是一个向上的步。那么相当于从 $(1,1)$ 出发向右向上走到 $(m+1,n)$
    \end{itemize}
\end{frame}
\begin{frame}{有限微积分}
    $$
    \begin{aligned}
    \sum_{1\le i_1\le \cdots\le i_m\le n}1
    &=\sum_{i_m=1}^n\sum_{i_{m-1}=1}^{i_m}\cdots\sum_{i_1=1}^{i_2}1=\sum_{i_m=1}^n\sum_{i_{m-1}=1}^{i_m}\cdots\sum_{i_2=1}^{i_3}i_2\\
    &=\sum_{i_m=1}^n\sum_{i_{m-1}=1}^{i_m}\cdots\sum_{i_3=1}^{i_4}\frac{i_3^{\overline{2}}}{2}\\
    &=\sum_{i_m=1}^n\sum_{i_{m-1}=1}^{i_m}\cdots\sum_{i_4=1}^{i_5}\frac{i_4^{\overline{3}}}{2\cdot 3}\\
    &\vdots\\
    &=\sum_{i_m=1}^n\frac{i_m^{\overline{m-1}}}{(m-1)!}=\frac{n^{\overline{m}}}{m!}=\binom{n+m-1}{m}
    \end{aligned}
    $$
\end{frame}
\begin{frame}{递推式}
    $$
    \begin{aligned}
    f_{m,n}&=\sum_{1\le i_1\le \cdots\le i_m\le n}1\\
    &=\sum_{i_m=1}^n\sum_{1\le i_1\le\cdots \le i_{m-1}\le i_m}\\
    &=\sum_{i_m=1}^n f_{m-1,n}
    \end{aligned}
    $$

    差分得到 $f_{m,n}-f_{m,n-1}=f_{m-1,n}$。
\end{frame}
\begin{frame}
    或者，在一开始就考虑
    $$
    \begin{aligned}
    f_{m,n}&=\sum_{1\le i_1\le \cdots\le i_m\le n}1\\
    &=\sum_{i_m=1}^n\sum_{1\le i_1\le\cdots \le i_{m-1}\le i_m<n}1+\sum_{i_m=1}^n\sum_{1\le i_1\le\cdots \le i_{m-1}\le i_m=n}\\
    &=f_{m,n-1}+f_{m-1,n}
    \end{aligned}
    $$
\end{frame}
\begin{frame}{解递推式}
    我们来考虑
    $$
    f_{n,m}=f_{n-1,m}+f_{n,m-1}+a_{n,m}
    $$
    这样的递推式，其中初值 $f_{\cdot,0},f_{0,\cdot}$ 任意指定。

    注意到 $a_{i,j}$ 对 $f_{n,m}$ 的贡献系数是：$g_{n,m}=g_{n-1,m}+g_{n,m-1}$ 在初值为 $g_{0,0}=1,g_{i,0}=g_{0,i}=0$ 时，$g_{n-i,m-j}$ 处的值。

    通过转组合意义，可以发现 $g_{n,m}$ 就是 $(0,0)$ 向右向上走到 $(n,m)$ 的方案数。所以 $a_{i,j}$ 对 $f_{n,m}$ 的贡献系数为 $\binom{n-i+m-j}{n-i}$。初值 $f_{i,0}$ 和 $f_{0,i}$ 的贡献系数同理。
\end{frame}
\begin{frame}{生成函数}
    TODO：我们在第五章再来处理这个问题
\end{frame}
\begin{frame}{一个 trick}
    $$
    \min(a_1,\cdots,a_n)=\sum_{i\ge 1}[\min(a_1,\cdots,a_n)\ge i]=\sum_{i\ge 1}[\forall k,\ a_k\ge i]
    $$

    $$
    \max(a_1,\cdots,a_n)=\sum_{i\ge 1}[\exists k,\ a_k\ge i]
    $$

    这样可以把最值问题转成判定问题来计数。

    另外，把 $f(i)$ 写成 $\sum_{1\le k\le i} f(k)-f(k-1)$ 常常是一个有效的思路（即使我们并没有在使用分部求和法！）。比如把 $x$ 写成 $\sum_i [1\le i\le x]$ 就是最常见的一个。

    把 $\left(\sum_i f(i)\right)^m$ 写成 $\sum_{i_1,i_2,\cdots,i_m}f(i_1)f(i_2)\cdots f(i_m)$ 常常是有用的。

    如果 $X$ 形如“满足xxx条件的元素的个数”，那么把 $X$ 写成 $\sum_x [x\text{满足条件}]$ 常常是有效的。这在\textbf{期望}中特别常见（称为“指示器随机变量”）。
\end{frame}
\begin{frame}{一种容斥}
    某些问题要求我们计数

    $$
    \sum (1-[P_1])(1-[P_2])\cdots (1-[P_k]).
    $$

    那么就可以拆括号变成

    $$
    \sum_{S\subseteq [k]} (-1)^{|S|}\sum \prod_{i\in S}[P_i],
    $$

    其中 $\sum\prod_i [P_i]$ 往往是容易算的。

    我觉得这个比那种标准的容斥要好理解和记忆得多。
\end{frame}
\begin{frame}{特殊情况}
    在前一页中，很多时候 $\sum\prod_{i\in S}[P_i]$ 只和 $|S|$ 有关。记这个和为 $N(|S|)$，此时有
    $$
    \begin{aligned}
    &\sum_{S\subseteq [k]} (-1)^{|S|}\sum \prod_{i\in S}[P_i]\\
    &=\sum_{S\subseteq [k]}(-1)^{|S|}N(|S|)\\
    &=\sum_i\binom{k}{i}(-1)^i N(i)
    \end{aligned}
    $$
\end{frame}
\begin{frame}{二项式反演}
    记 $g_i=\sum_{|S|=i}\sum\prod_{i\in S}[P_i]$ 为“钦定了 $i$ 个性质的方案数”，那么就有
    $$
    \begin{aligned}
    f_k&=\text{恰好满足}\ k\text{ 个性质的方案数}=\sum_{|S|=k}\sum \prod_{i\in S}[P_i]\prod_{j\notin S}(1-[P_j])\\
    &=\sum_{|S|=k}\sum\left(\prod_{i\in S}[P_i]\right)\sum_{T\subseteq \bar S}(-1)^{|T|}\prod_{j\in T}[P_j]\\
    &=\sum_{|S|=k}\sum_{T'\supseteq S}(-1)^{|T'|-|S|}\sum\prod_{i\in T'}[P_i]\qquad(T'=S\uplus T)\\
    &=\sum_{n}\left(\sum_{|T'|=n}\sum\prod_{i\in T'}[P_i]\right)\left(\sum_{S\subseteq T',|S|=k}(-1)^{|T'|-|S|}\right)\\
    &=\sum_n g_n\binom{n}{k}(-1)^{n-k}
    \end{aligned}
    $$
\end{frame}
\begin{frame}{二项式反演（续）}
    练习：反过来用 $f$ 表示 $g$。
    
    设 $h_i=\sum_{|S|=i}\sum\prod_{i\notin S}[\overline{P_i}]$ 为“至多满足 $i$ 个性质的方案数”，找出 $h$ 和 $f$ 之间的关系。
\end{frame}
\begin{frame}{分部求和法：另一种视角}
    $$
    \begin{aligned}
    \sum_{i=1}^na_ib_i&=\sum_{i=1}^na_i\left(b_0+\sum_{j=1}^i(b_j-b_{j-1})\right)\\
    &=b_0\sum_{i=1}^na_i+\sum_{j=1}^n(b_j-b_{j-1})\sum_{i=j}^n a_i\\
    &=b_0\sum_{i=1}^na_i+\sum_{j=1}^n(b_j-b_{j-1})\left(\sum_{i=1}^na_i-\sum_{i=1}^{j-1}a_i\right)\\
    &=b_n\sum_{i=1}^na_i-\sum_{j=1}^n(b_j-b_{j-1})\sum_{i=1}^{j-1}a_i
    \end{aligned}
    $$
    （注意以上的第二行和第四行都是可用的结论）

    也就是说，对 $a_ib_i$ 求和可以转为对 $(\sum a_i)(\Delta b_i)$ 求和
\end{frame}
\begin{frame}{例子}
    给一棵树，点有点权，求所有路径 $\min$ 的 $k$ 次方的和。（hint：先考虑 $k=1$）

    $1,\ldots,n$ 的所有排列的逆序对数之和是多少？

    长为 $n$ 的序列，每个都是 $1$ 到 $m$ 中的数。一个序列的价值为其中不同数的数量，求所有序列的价值和。（有封闭形式）

    长为 $n$ 的序列，每个都是 $1$ 到 $m$ 中的数。一个序列的价值为不同数的数量的 $k$ 次方，求所有序列的价值和。（你可以认为 $k$ 非常小或者干脆就是个常数，虽然我们实际上至少可以做到 $O(k(\log k+\log n))$） 

    长为 $n$ 的序列，每个都是 $1$ 到 $m$ 中的数。要求 $1,\cdots,k$ 必须在序列中出现，求方案数。（事实上这是上一个问题的 $O(k\log k)$ 做法的一个子问题，你只需要给出一个 $O(k^2)$ 的做法）
\end{frame}
\begin{frame}{路径 $\min$ 的 $k$ 次方和}
    $$
    \begin{aligned}
        &\sum_{\text{路径}P}\min(P)^k=\sum_{\text{路径}P}\sum_{i=1}^{\min(P)}i^k-(i-1)^k\\
        &=\sum_{i\ge 1}(i^k-(i-1)^k)\sum_{\text{路径}P}[\min(P)\ge i]\\
        &=\sum_{i\ge 1}(i^k-(i-1)^k)\sum_{\text{路径}P}[P\text{ 是只保留权值 }\ge i \text{ 的点形成的森林中的一条路径}]\\
        &=\sum_{i\ge 1}(i^k-(i-1)^k)\sum_{B}[B\text{ 是只保留权值 }\ge i\text{ 的点形成的森林中的一个连通块}]\binom{|B|+1}{2}
    \end{aligned}
    $$
    假设点权从大到小为 $w_m,\ldots,w_1$。在 $i\in [w_i,w_{i+1})$ 时第二个求和的结果都相同，而 $i$ 从 $w_i$ 减小到 $w_i-1$ 时连通块会发生合并。并查集维护连通块大小即可。
\end{frame}
\begin{frame}{路径 $\min$ 的 $k$ 次方和：另解}
    也可以直接枚举 $\min$ 的取值：
    $$
    \begin{aligned}
        &\sum_{\text{路径}P}\min(P)^k=\sum_w w^k\sum_{\text{路径}P}[\min(P)=w]\\
        &=\sum_{w}w^k\sum_{\text{路径}P}([\min(P)\ge w]-[\min(P)>w]\\
    \end{aligned}
    $$

    剩下的过程是一样的。通过分部求和法也不难发现这两个式子相等。
\end{frame}
\begin{frame}{排列逆序对数之和}
    $$
    \begin{aligned}
    \sum_{\pi}\mathrm{inv}(\pi)&=\sum_{\pi}\sum_{i<j}[\pi_i>\pi_j]\\
    &=\sum_{i<j}\sum_{\pi}[\pi_i>\pi_j]\\
    &=\sum_{i<j}\frac{n!}{2}\\
    &=\frac{n!}{2}\binom{n}{2}
    \end{aligned}
    $$

    其中第二行到第三行使用了排列的对称性，即 $\pi_i>\pi_j$ 和 $\pi_i<\pi_j$ 的排列个数是一样多的。
\end{frame}
\begin{frame}{序列价值和 I}
    $$
    \begin{aligned}
    \sum_{a}w(a)&=\sum_a\sum_{i=1}^m[i\text{ 在 }a\text{ 中出现}]\\
    &=\sum_{i=1}^m\sum_a[i\text{ 在 }\ a\text{ 中出现}]\\
    &=\sum_{i=1}^m\sum_a(1-[i\text{ 不在 }\ a\text{ 中出现}])\\
    &=\sum_{i=1}^m(m^n-(m-1)^n)\\
    &=m(m^n-(m-1)^n)
    \end{aligned}
    $$
\end{frame}
\begin{frame}{序列价值和 I：另解}
    $$
    \begin{aligned}
    \sum_{a}w(a)&=\sum_a\sum_{i=1}^n[a_i\text{ 这个值在 }\ a\text{ 中首次出现}]\\
    &=\sum_a\sum_{i=1}^n[a_i\text{ 这个值在 }\ \{a_1,\ldots,a_{i-1}\}\text{ 中没有出现}]\\
    &=\sum_{i=1}^n\sum_{x=1}^m\sum_{a}[a_i=x][x\text{ 在 }\ \{a_1,\ldots,a_{i-1}\}\text{ 中没有出现}]\\
    &=\sum_{i=1}^n\sum_{x=1}^m(m-1)^{i-1}m^{n-i}=\sum_{i=1}^n(m-1)^{i-1}m^{n-i+1}\\
    &=m^{n}\frac{1-((m-1)/m)^n}{1-(m-1)/m}=m(m^n-(m-1)^n)
    \end{aligned}
    $$
\end{frame}
\begin{frame}{序列价值和 II}
    $$
    \begin{aligned}
    \sum_{a}w(a)^k&=\sum_a\sum_{1\le i_1,i_2,\ldots,i_k\le m}[i_1,\ldots,i_k\text{ 都在 }\ a\text{ 中出现}]\\
    &=\sum_{1\le i_1,i_2,\ldots,i_k\le m}\sum_a[i_1,\ldots,i_k\text{ 都在 }\ a\text{ 中出现}]
    \end{aligned}
    $$

    观察小的情况，可以发现
    $$
    \sum_a[i_1,\ldots,i_k\text{ 都在 }\ a\text{ 中出现}]
    $$
    只与 $i_1,\ldots,i_k$ 中不同值的个数相关。并且，假设恰有 $t$ 个不同值，那么由对称性有
    $$
    \sum_a[i_1,\ldots,i_k\text{ 都在 }\ a\text{ 中出现}]=\sum_a[1,\ldots,t\text{ 都在 }\ a\text{ 中出现}]
    $$
\end{frame}
\begin{frame}
    枚举 $i_1,\ldots,i_k$ 中不同值的个数：
    $$
    \begin{aligned}
        &\sum_{1\le i_1,i_2,\ldots,i_k\le m}\sum_a[i_1,\ldots,i_k\text{ 都在 }\ a\text{ 中出现}]\\
        &=\sum_t\sum_{1\le i_1,i_2,\ldots,i_k\le m}[i_1,\ldots,i_k\text{ 恰有 }\ t\text{ 个不同值}]\sum_a[i_1,\ldots,i_k\text{ 都在 }\ a\text{ 中出现}]\\
        &=\sum_t\left(\sum_{1\le i_1,i_2,\ldots,i_k\le m}[i_1,\ldots,i_k\text{ 恰有}\ t\text{ 个不同值}]\right)\left(\sum_a[1,\ldots,t\text{ 都在}\ a\text{ 中出现}]\right)
        \end{aligned}
    $$
\end{frame}
\begin{frame}
    前面这部分是
    $$
    \sum_{1\le x_1<x_2<\cdots<x_t\le m}\sum_{1\le i_1,i_2,\ldots,i_k\le m}[\{i_1,\ldots,i_k\}\text{ 这个集合等于 }\{x_1,\ldots,x_t\}].
    $$
    由对称性，
    $$
    \sum_{1\le i_1,i_2,\ldots,i_k\le m}[\{i_1,\ldots,i_k\}=\{x_1,\ldots,x_t\}]=\sum_{1\le i_1,i_2,\ldots,i_k\le m}[\{i_1,\ldots,i_k\}=\{1,\ldots,t\}].
    $$
    接下来
    $$
    \begin{aligned}
    &\sum_{1\le x_1<x_2<\cdots<x_t\le m}\sum_{1\le i_1,i_2,\ldots,i_k\le m}[\{i_1,\ldots,i_k\}\text{ 这个集合等于 }\{1,\ldots,t\}]\\
    &=\binom{m}{t}\sum_{1\le i_1,i_2,\ldots,i_k\le t}[1,\ldots,t\text{ 都在 }(i_1,\ldots,i_k)\text{ 中出现}].
    \end{aligned}
    $$
    注意到这和后半部分是相同形式的问题。
\end{frame}
\begin{frame}
    令 
    $$
    f(n,m,k)=\sum_{1\le i_1,\ldots,i_n\le m}[1,\ldots,k\text{ 都在 }(i_1,\ldots,i_n)\text{ 中出现}].
    $$
    容斥：
    $$
    \begin{aligned}
    f(n,m,k)&=\sum_{1\le i_1,\ldots,i_n\le m}\prod_{j=1}^k(1-[j\text{ 在 }\{i_1,\ldots,i_n\}\text{ 中不出现}])\\
    &=\sum_{1\le i_1,\ldots,i_n\le m}\sum_{S\subseteq \{1,\ldots,k\}}(-1)^{|S|}[\forall x\in S,\ x\text{ 在 }\{i_1,\ldots,i_n\}\text{ 中不出现}]\\
    &=\sum_{S\subseteq \{1,\ldots,k\}}(-1)^{|S|}\sum_{1\le i_1,\ldots,i_n\le m}[\forall x\in S,\ x\text{ 在 }\{i_1,\ldots,i_n\}\text{ 中不出现}]\\
    &=\sum_{S\subseteq \{1,\ldots,k\}}(-1)^{|S|}(m-|S|)^n\\
    &=\sum_t\binom{k}{t}(-1)^{t}(m-t)^n
    \end{aligned}
    $$
\end{frame}
\begin{frame}
    现在
    $$
    \begin{aligned}
        &\sum_t\left(\sum_{1\le i_1,i_2,\ldots,i_k\le m}[i_1,\ldots,i_k\text{ 恰有}\ t\text{ 个不同值}]\right)\left(\sum_a[1,\ldots,t\text{ 都在}\ a\text{ 中出现}]\right)\\
        &=\sum_{t=0}^k f(k,t,t)f(n,m,t)
    \end{aligned}
    $$
    而
    $$
    f(n,m,k)=\sum_t\binom{k}{t}(-1)^t(m-t)^n
    $$
    可以 $O(k)$ 计算。故本题可以 $O(k^2)$ 完成。实际上，通过卷积也可以 $O(k\log k)$ 完成。
    
    事实上，$f(n,k,k)$ 就是第二类斯特林数 $k! {n \brace k}$：把 $1,\ldots,k$ 视作有标号盒子，$i_1,\ldots,i_n$ 视作球即可。
\end{frame}
\begin{frame}{序列价值和 II：另解}
    如果在一开始就枚举序列的价值，就得到
    $$
    \begin{aligned}
    \sum_w w^k\sum_a [w(a)=w]&=\sum_w w^k\sum_a [a\text{ 中恰有}\ w\text{ 个不同值}]\\
    &=\sum_w w^k\binom{m}{w}\sum_a[a\text{ 中出现的元素恰好是}\ \{1,\ldots,w\}]\\
    &=\sum_w w^k\binom{m}{w}f(n,w,w)
    \end{aligned}
    $$
    但是这个并不能直接用于计算，因为 $w$ 有意义的范围高达 $[0,\min(n,m)]$。但是我们可以用前面得到的 $f$ 的公式继续化简。TODO：我们在学习第五章之后再来处理这个式子。
\end{frame}
% \begin{frame}
%     $$
%     \begin{aligned}
%       \sum_w w^k\binom{m}{w}f(n,w,w)&=\sum_w w^k\binom{m}{w}\sum_t \binom{w}{t}(-1)^t(w-t)^n\\
%       &=\sum_t(-1)^t\binom{m}{t}\sum_w w^k\binom{m-t}{w-t}(w-t)^n\\
%       &=\sum_t(-1)^t\binom{m}{t}\sum_w (w+t)^k\binom{m-t}{w}w^n\\
%       &=\sum_t (-1)^t\binom{m}{t}\sum_w\binom{m-t}{w}w^n\sum_j\binom{k}{j}w^jt^{k-j}\\
%       &=\sum_t (-1)^t\binom{m}{t}\sum_w\binom{m}{w}(w-t)^n\sum_j{k\brace j}j!\binom{w}{j}\\
%       &=\sum_t (-1)^t\sum_j{k\brace j}j!\binom{m}{j}\sum_w\binom{m-j}{w-j}(w-t)^n  
%     \end{aligned}
%     $$
% \end{frame}
% \begin{frame}
%     $$
%     \begin{aligned}
%     \sum_n\frac{z^n}{n!}\sum_w w^k\binom{m}{w}{n\brace w}w!&=\sum_w\binom{m}{w}w^kw!\sum_n\frac{z^n}{n!}{n\brace w}\\
%     &=\sum_w\binom{m}{w}w^k(e^z-1)^w\\
%     &=\sum_w\binom{m}{w}(e^z-1)^w\sum_j{k\brace j}j!\binom{w}{j}\\
%     &=\sum_j{k\brace j}j!\binom{m}{j}\sum_w(e^z-1)^{w-j+j}\binom{m-j}{w-j}\\
%     &=\sum_j{k\brace j}j!\binom{m}{j}(e^z-1)^je^{(m-j)z}\\
%     \end{aligned}
%     $$
% \end{frame}
\begin{frame}{例子：变式}
    给一个 $0$ 到 $n-1$ 的排列，求所有区间的 $\mathrm{mex}$ 的 $k$ 次方的和。$\mathrm{mex}(S)$ 为 $S$ 中最小的没有出现过的自然数。（Bonus：排列区间->序列区间；排列区间->树路径（权值为排列）；树路径（不假设权值为排列能做吗？或者证明 hardness））

    $1,\ldots,n$ 的所有排列 $\pi$ 的\textbf{严格前缀} $\mathbf{\max}$ \textbf{数量} 之和是多少？一个位置 $i$ 是严格前缀 $\max$ 当且仅当 $\forall j<i.\pi_i>\pi_j$。

    长为 $n$ 的序列，每个都是 $1$ 到 $m$ 中的数。定义一个序列 $a$ 的权值 $w(a)$ 为 $a$ 中不同数的数量。完成以下两个问题中的至少一个（需要 $O(k^2)$）：
    \begin{itemize}
        \item 求所有序列的 $(w(a))^{\underline{k}}$ 的和
        \item 给定 $k$ 次多项式 $f$，求所有序列的 $f(w(a))$ 的和
    \end{itemize}
\end{frame}

\begin{frame}{[CSP-S 2019 江西] 和积和}
    $$
    \sum_{l=1}^n\sum_{r=1}^n\sum_{i=l}^ra_i\sum_{i=l}^r b_i.
    $$

    需要 $O(n)$
\end{frame}
\begin{frame}{数列求和}
    $$
    \sum_{i=1}^n i^ka^i
    $$

    需要 $O(k^2)$.

    分别使用扰动法和分部求和法完成
\end{frame}
\begin{frame}
    设 $S_n=\sum_{i=1}^n i^k a^i$，那么首先
    $$
    S_n=\sum_{i=1}^n (i+1)^k a^{i+1} - (n+1)^ka^{n+1}+1^ka^{1}.
    $$

    扰动法的关键一步在于把 $(i+1)^ka^{i+1}$ 表示成关于 $i$ 而不是 $i+1$ 的表达式。在遇到和的幂时，我们的第一反应就是二项式定理：
    $$
    (i+1)^ka^{i+1}=a\sum_j\binom{k}{j}i^ja^i.
    $$ 

    把这个带到扰动法中就得到
    $$
    S_n=\sum_{i=1}^n a\sum_j\binom{k}{j}i^ja^i=a\sum_j\binom{k}{j}\sum_{i=1}^n i^ja^i. 
    $$
    
    （要牢记我们的基本准则“提出无关项”和“交换求和号”）。
\end{frame}
\begin{frame}
    现在，我们发现内部的求和和 $S_n$ 很像，只是 $k$ 换成了 $j$。这种时候我们往往可以尝试设出辅助递归项 $S_{n,j}$，而 $S_{n,j}$ 关于 $j$ 的递归具有同样的模式。

    特别地，在 $a=1$ 时无法得到 $S_{n,k}$ 的方程，而是 $S_{n,k-1}$ 的方程。然而，这并不影响我们能够递推求解 $S_{n,j}$ 的事实。

    这样我们就得到了 $O(k^2)$ 的做法。而如果我们挥动线性递推重锤，我们可以得到一个 $O(k)$ 的做法。具体来说：
    \begin{itemize}
        \item 线性递推告诉我们结果具有 $a^nP(n)+C$ 的形式，其中 $P(n)$ 是一个 $k$ 次多项式，$C$ 是常数
        \item 先递推计算出 $S_0,\ldots,S_k$，然后 $P(i)=a^{-i}(S_i-C)$
        \item 使用 Lagrange 插值和直接递推两种方式分别计算 $P(k+1)$，从而解出 $C$
        \item 使用 Lagrange 插值得到 $P(n)$ 和 $S_n$
    \end{itemize}
\end{frame}
\begin{frame}{分部求和法}
    自行尝试
\end{frame}
\begin{frame}{普通幂和下降幂的互相转换}
    我们可以在 $O(k^2)$（甚至 $O(k\log^2 k)$）内在普通多项式和下降幂多项式之间相互转换。

    做法一：拉格朗日插值

    做法二：可以在 $O(k)$ 内从 $x^{\underline{k-1}}$ 计算 $x^{\underline{k}}$，从而 $O(k^2)$ 计算所有 $x^{\underline{1}},\cdots,x^{\underline{k}}$。反过来，可以类似多项式除法从高到低计算 $\sum_{i}a_ix^i$ 转到 $\sum_i b_ix^{\underline{i}}$ 后的各项系数。

    做法三：
    $$
    x^k=\sum_i \begin{Bmatrix}k\\i\end{Bmatrix}x^{\underline{i}},\quad x^{\underline{k}}=\sum_i\begin{bmatrix}k\\i\end{bmatrix}x^i(-1)^{k-i}
    $$
\end{frame}
\begin{frame}{斯特林数}
    我如何在不知道斯特林数的定义的情况下计算斯特林数？我们有 $\displaystyle  x^k=\sum_i \begin{Bmatrix}k\\i\end{Bmatrix}x^{\underline{i}}$，而另一方面有
    $$
    \begin{aligned}
    x^k&=x\cdot x^{k-1}=\sum_i\begin{Bmatrix}k-1\\i\end{Bmatrix}x\cdot x^{\underline{i}}=\sum_i\begin{Bmatrix}k-1\\i\end{Bmatrix}\cdot x^{\underline{i}}(x-i+i)\\
    &=\sum_i\begin{Bmatrix}k-1\\i\end{Bmatrix}\cdot (ix^{\underline{i}}+x^{\underline{i+1}})\\
    &=\sum_i\left(\begin{Bmatrix}k-1\\i-1\end{Bmatrix}+i\begin{Bmatrix}k-1\\i\end{Bmatrix}\right)x^{\underline{i}}.
    \end{aligned}
    $$

    比较系数就得到 $\displaystyle\begin{Bmatrix}k\\i\end{Bmatrix}=\begin{Bmatrix}k-1\\i-1\end{Bmatrix}+i\begin{Bmatrix}k-1\\i\end{Bmatrix}$，另一个同理
\end{frame}
\begin{frame}{二项式定理（负指数）}
    计算
    $$
    \sum_{i=0}^ni^{\overline{k}}r^i
    $$

    假设 $0<r<1$，令 $n=\infty$，然后重新审视你的结果
\end{frame}
\begin{frame}{三角形：直接的推法 1}
    $$
    \begin{aligned}
        &\sum_{i,j,k}[1\le i\le j\le k\le n][i+j+k=n][i+j>k]\\
        &=\sum_{i,j}[1\le i\le j\le n-i-j][i+j>n-i-j]\\
        &=\sum_{1\le i\le j}[i\le n-2j][i>\frac{n}{2}-j]\\
        &=\sum_{j\ge 1}\sum_{i\ge 1}[\frac{n}{2}-j<i\le \min(j,n-2j)]
    \end{aligned}
    $$

    接下来分类讨论 $\min(j,n-2j)$ 是哪个。当 $j\le n-2j$ 也就是 $j\le n/3$ 时是 $j$，否则也就是 $j>n/3$ 时是 $n-2j$。
\end{frame}
\begin{frame}
    $j\le n/3$ 的部分如下：
    $$
    \sum_{1\le j\le n/3}\sum_{i\ge 1}[\frac{n}{2}-j<i\le j].
    $$
    在 $i$ 的合法区间非空，也就是 $j<n/4$ 时，合法的 $i$ 有 $j-\lfloor n/2-j\rfloor=2j-\lfloor n/2\rfloor$ 个。所以这部分是
    $$
    \sum_{n/4<j\le n/3}(2j-\lfloor n/2\rfloor),
    $$
    等差数列求和。
\end{frame}
\begin{frame}
    $j>n/3$ 的部分如下：
    $$
    \sum_{n/3<j}\sum_{i\ge 1}[\frac{n}{2}-j<i\le n-2j].
    $$
    在 $i$ 的合法区间非空，也就是 $j<n/2$ 时，合法的 $i$ 有 $n-2j-\lfloor n/2-j\rfloor=\lceil n/2\rceil-j$ 个。所以这部分是
    $$
    \sum_{n/3<j<n/2}(\lceil n/2\rceil-j),
    $$
    同样等差数列求和。
\end{frame}
\begin{frame}{三角形：直接的推法 2}
    在最开始的条件转化中，除了先对 $j$ 再对 $i$ 求和外，也可能会变成先对 $i$ 求和再对 $j$ 求和。注意到这样子势必会引入 $j\le\frac{n-i}{2}$ 这样的条件，从而内层对 $j$ 的求和会带上 $\lfloor (n-i)/2\rfloor$ 这样的东西。回顾我们的第一条化简准则——保持求和的条件简单，这样的转化不应该首先被考虑。首先考虑的应该是计算不等式中不带系数的 $i$ 的合法区间。
    
    不过，即使出现了 $\lfloor i/2\rfloor$ 这样的求和也不要担心，我们可以将其写成 $\frac{i-(i\bmod 2)}{2}$ 的形式，而 $i\bmod 2=[i\text{ 是奇数}]$。所以在计算 $\sum_{i} T(i)(i\bmod 2)$ 时，可以做代换 $i=2k+1$ 变成 $\sum_{2k+1}T(2k+1)$。

    我们将在第三章见到更多更一般的处理带有取整的求和的方法。
\end{frame}
\begin{frame}{三角形：容斥的推法}
    使用容斥可以得到简单一些的推法。
    $$
    \begin{aligned}
        &\sum_{i,j,k}[1\le i\le j\le k\le n][i+j+k=n][i+j>k]\\
        &=\sum_{i,j,k}[i+j+k=n][1\le i\le j\le n-i-j\le n](1-[i+j\le n-i-j])\\
        &=\sum_{1\le i\le j}[j\le n-i-j]-\sum_{1\le i\le j}[j\le n-i-j][i+j\le n-i-j]
    \end{aligned}
    $$
\end{frame}
\begin{frame}
    前半部分是
    $$
    \begin{aligned}
        &\sum_{1\le i\le j}[j\le n-i-j]\\
        &=\sum_{1\le j}\sum_{i=1}^j[i\le n-2j]\\
        &=\sum_{1\le j\le n/3}\sum_{i=1}^j 1+\sum_{n/3<j}\sum_{i=1}^{n-2j}1\\
        &=\sum_{1\le j\le \lfloor n/3\rfloor}j+\sum_{j=\lfloor n/3\rfloor+1}^{\lfloor n/2\rfloor}(n-2j)\\
    \end{aligned}
    $$    

    等差数列求和。
\end{frame}
\begin{frame}
    后半部分是
    $$
    \begin{aligned}
        &\sum_{1\le i\le j}[j\le n-i-j][i+j\le n-i-j]=\sum_{1\le i\le j}[i+j\le n-i-j]\\
        &=\frac{1}{2}\left(\sum_{1\le i,j}[i+j\le n-i-j]+\sum_{1\le i=j}[i+j\le n-i-j]\right)\\
        &=\frac{1}{2}\left(\sum_{1\le i,j}\sum_d[i+j=d][i+j\le n/2]+\sum_{1\le i=j}[i+j\le n/2]\right)\\
        &=\frac{1}{2}\left(\sum_d\sum_{1\le i,j}[i+j=d][d\le n/2]+\sum_{1\le i}[i\le n/4]\right)\\
        &=\frac{1}{2}\left(\sum_{d=2}^{\lfloor n/2\rfloor}(d-1)+\lfloor n/4\rfloor\right)\\
    \end{aligned}
    $$
    同样是等差数列求和    
\end{frame}
\begin{frame}{三角形：对称性的推法}
    首先，注意到 $P(i,j,k)=[i+j+k=n][i,j,k\text{ 构成三角形}]$ 是完全对称的。如果你知道 Burnside 引理，或者尝试使用对称性并手算系数的话，可以发现
    $$
    \sum_{1\le i\le j\le k}P(i,j,k)=\frac{1}{6}\left(2\sum_{1\le i=j=k} P(i,j,k)+3\sum_{1\le i=j,k}P(i,j,k)+\sum_{1\le i,j,k}P(i,j,k)\right).
    $$
    接下来，注意到
    $$
    \begin{aligned}
    P(i,j,k)&=[i+j+k=n][i,j,k\text{ 构成三角形}]\\
    &=[i+j+k=n][i+j+k>2\max(i,j,k)]\\
    &=[i+j+k=n][\max(i,j,k)<\frac{n}{2}]\\
    &=[i+j+k=n][i,j,k<\frac{n}{2}]\\
    &=[i+j+k=n](1-[i\ge n/2])(1-[j\ge n/2])(1-[k\ge n/2])\\
    \end{aligned}
    $$
\end{frame}
\begin{frame}
    注意在 $1\le i,j,k$ 的限制下，拆括号后如果一项中出现了两个或更多 $[i\ge n/2]$ 的乘积，那么这个项不会有贡献（因为此时 $i+j+k$ 必定超过 $n$）。类似地，$i=j,k$ 这个求和里面不会拆出来 $[i\ge n/2]$ 的项。所以我们有
    $$
    \sum_{1\le i=j=k}P(i,j,k)=\sum_{1\le i}[3i=n]=[3\mid n],
    $$
    $$
    \begin{aligned}
    \sum_{1\le i=j,k}P(i,j,k)&=\sum_{1\le i=j,k}[i+j+k=n](1-[k\ge n/2])\\
    &=\sum_{1\le i}[n-2i\ge 1]-\sum_{1\le i}[n/2\le n-2i]\\
    &=\left\lfloor\frac{n-1}{2}\right\rfloor-\left\lfloor\frac{n}{4}\right\rfloor.
    \end{aligned}
    $$
\end{frame}
\begin{frame}
    以及
    $$
    \begin{aligned}
    &\sum_{1\le i,j,k}P(i,j,k)\\
    &=\sum_{1\le i,j,k}[i+j+k=n](1-[i\ge n/2]-[j\ge n/2]-[k\ge n/2])\\
    &=\sum_{1\le i,j,k}[i+j+k=n]-3\sum_{1\le i,j,k}[i+j+k=n][i\ge n/2]\\
    &=\sum_{1\le i,j,k}[i+j+k=n]-3\sum_{1\le (i-\lceil n/2\rceil+1),j,k}[(i-\lceil n/2\rceil+1)+j+k=n-(\lceil n/2\rceil-1)]\\
    &=\binom{n-1}{2}-3\binom{\lfloor n/2\rfloor-2}{2},
    \end{aligned}
    $$

    这个做法可以比较方便地从三角形扩展到一般的 $k$ 边形（在旋转和翻转下数等价类）
\end{frame}
